% =====================================================================
% SEC 01 — Introdução
% =====================================================================
\section{Introdução}
\label{sec:intro}

Sistemas de \emph{Retrieval-Augmented Generation} (RAG) combinam recuperação de
documentos com geração de linguagem natural para responder consultas com
\emph{grounding} em evidências \citep{Lewis2020RAG}. Em domínios científicos, a
confiabilidade da resposta depende diretamente da qualidade e da \emph{cobertura}
do conjunto de documentos recuperados \citep{Ram2023RetrievalAugmented}. Contudo,
os estágios de ranqueamento de tais sistemas são, em sua grande maioria,
orientados a \emph{relevância} individual \citep{RobertsonZaragoza2009BM25},
medindo quão bem cada documento responde à consulta isoladamente — e não quão
\emph{diversos} ou \emph{complementares} eles são entre si.

\begin{quote}\small
\textbf{Ponto-chave.} Para consultas de \emph{paisagem} ou \emph{multicunha} (revisão de literatura,
síntese de evidências, comparação entre perspectivas), um conjunto recuperado
que concentra-se em um único ângulo subestima o espaço de respostas mesmo quando
cada item individualmente é altamente relevante. A diversidade dos candidatos é,
nesses casos, tão importante quanto a relevância individual.
\end{quote}

Do ponto de vista de mensuração, arquiteturas científicas de RAG frequentemente
expõem um painel de métricas de \emph{groundedness}, cobertura de citação e
propagação temporal, mas \textbf{ausentam} uma métrica de diversidade do conjunto
recuperado. Essa lacuna metodológica motivou a formulação da \emph{proposta
pós-Recamán} \citep{Alekseyev2022RecamanCousins}: um estágio de diversificação
\emph{determinístico} e \emph{pós-ranqueado}, baseado na sequência de Recamán
(OEIS A005132), que gera \emph{offsets} fixos sobre o ranking para espalhar a
seleção, acompanhado de uma métrica de diversidade $\Div(\mathcal{S})$ sobre
\emph{âncoras canônicas} de conteúdo.

A proposta é atraente por três razões: (i) \emph{determinismo} — os offsets são
fixos e universais, eliminando a aleatoriedade da amostragem; (ii) \emph{custo
baixo} — geração da sequência em $O(M)$ e seleção em $O(K)$ sobre o ranking; e
(iii) \emph{interpretabilidade} — não depende de parâmetros como o $\lambda$ do
MMR \citep{CarbonellGoldstein1998MMR}. A questão que este artigo investiga é,
porém, essencialmente \emph{empírica}: dado um ranking, a seleção posicional por
Recamán de fato aumenta a diversidade de \emph{conteúdo} (âncora/fonte) quando
comparada ao \emph{top-k}? E como ela se compara ao MMR?

Para respondê-la de forma controlada e reprodutível, construímos um
\emph{benchmark} de coorte com corpus-piloto determinístico, no qual famílias de
documentos representam ângulos de perspectiva distintos, e comparamos quatro
estratégias de seleção pós-ranqueamento sob métricas de diversidade, relevância
e cobertura. O resultado central é \emph{factual} e \emph{contrarintuitivo} em
relação à expectativa inicial: no cenário testado, o esquema posicional de
Recamán \textbf{empata} com o \emph{top-k} em diversidade; o MMR supera ambos;
e um novo diversificador híbrido, \textbf{HABD}, que combina determinismo com
dissimilaridade de \emph{âncoras} e $\lambda$ adaptativo por query, supera ambos
em diversidade e cobertura — porém ao custo de uma queda de relevância acima da
tolerância adotada. Esses achados evidenciam (i) uma \emph{limitação estrutural}
de diversificadores posicionais e (ii) que \emph{$\lambda$ fixo é subótimo},
alinhando-se à literatura recente de \emph{diversidade ciente da consulta}
\citep{Khan2026DFRAG,RezaeiDieng2025VendiRAG}.

\subsection{Contribuições}

Resumimos as contribuições deste artigo:

\begin{enumerate}
  \item \textbf{C1 (formalização)} — a formulação matemática da diversificação
        pós-ranqueamento por Recamán e da métrica $\Div(\mathcal{S})$ sobre
        âncoras canônicas, de caráter aditivo e determinístico;
  \item \textbf{C2 (método)} — um \emph{benchmark} de coorte aberto e
        reprodutível para comparar estratégias de diversificação
        (\texttt{benchmarks/cohort\_recaman.py}), incluindo \emph{top-k}, MMR,
        Recamán e HABD;
  \item \textbf{C3 (evidência)} — a evidência empírica controlada de que
        diversificadores posicionais podem coincidir com o \emph{top-k} sob
        ranking monopolista por família, enquanto o MMR, baseado em
        dissimilaridade de conteúdo, espalha por fonte;
  \item \textbf{C4 (HABD)} — um diversificador híbrido \emph{anchor-blended}
        determinístico com $\lambda$ adaptativo por query (sem LLM), que
        maximiza a cobertura de âncoras; e a avaliação honesta de que, apesar
        de superar MMR/Recamán em diversidade, \textbf{não satisfaz} a
        tolerância de relevância (H3 refutada) — reporte equilibrado raro na
        literatura.
\end{enumerate}

\subsection{Organização do artigo}

Na Seção~\ref{sec:rel} revisamos trabalhos relacionados, incluindo a literatura
recente de diversidade \emph{adaptativa} em RAG (DF-RAG, Vendi-RAG, PLMMR). A
Seção~\ref{sec:metodo} apresenta a formulação das propostas (Recamán e HABD), o
desenho do corpus-piloto e o protocolo de avaliação. A Seção~\ref{sec:resultados}
reporta os resultados observados para as quatro estratégias. A
Seção~\ref{sec:discussao} discute implicações e limitações, e a
Seção~\ref{sec:conclusao} conclui. Declaramos, desde já, o escopo: todos os
valores quantitativos provêm de \emph{corpus piloto controlado} e \emph{não
generalizam para produção}.
